\scanpage{120}
$X$ schéma projectif
\[
M\subset K(X)
\]
\uline{Conditions équivalentes} \quad (i) $\exists$ familles \uline{limitées} de faisceaux \inserted{loc. libres} $(F_i)(G_i)$ sur $X$, tel que $M$ contenu ds l'ens des
\[
\gamma(F_i)-\gamma(G_i)
\]

(ii) $\exists$ un $n,N$ et des \unsure{partie|parties} finie \deleted{de polynômes \unsure{en}} \deleted{$\unsure{P_1,Q_i|P_1,P_i}$} $\unsure{\mathcal P|\Phi|P}\subset Q[T]$, telle que les $x\in M$ soient de la forme $\gamma(F)-\gamma(G)$, $F$ et $G$ des quotients loc. libres de $\mathcal O(-n)^N$, de polynômes de Hilbert $\in\mathcal P$.

(iii) \inserted{Les polynômes de Hilbert des $\xi\in M$ en nb fini, et} $\exists$ \deleted{des} faisceaux loc. libres $\deleted{\unsure{F}}A$ et $B$ ($\underline{\mathcal O}(-n)^N$ si on veut) tels que les $\xi+\gamma(A)$ $(\xi\in M)$ soient décrits par des \uline{faisceaux} loc libres, quotients de \deleted{un $=$ faisceau $\underline{\mathcal O}$ loc. libre \unsure{($\underline{\mathcal O}(-n)^N$|($\mathcal O(-n)^N$} si on veut} \inserted{$B$.} $]$.

(iii) $\Longrightarrow$ (ii) trivial, compte tenu que $A,B$ sont quotients de $\underline{\mathcal O}(-n)^N$ pour $n,N$ convenant,

(ii) $\Longleftrightarrow$ (i) \quad Par la caractérisation des familles limitées

(ii) $\Longrightarrow$ (iii) \quad Car si $G=O(-n)^N/R$, on \unsure{aura}
\[
\xi=\gamma(F)-\gamma(\underline O(-n)^N)+\gamma(R)\qquad\text{d'où}
\]
\[
\xi+\gamma(\underline{\mathcal O}(-n)^N)=\gamma(F+R).
\]
Or $\unsure{F|E}+R$ est un quotient de $\underline O(-\unsure{n|1})^{N}\unsure{'|\}|_{\xi}}$, donc les $F+\text{\uuline{$R$}}$ \uline{sont quotients} d'un

% archive sheet number: 119
